% 第九卷 13.1.3 节：检索增强生成流水线及各改进方法作用的环节（跨栏，约152 mm宽）。
% 实线：一次检索增强生成的数据流；虚线：由生成过程触发的再检索。
% 蓝色：输入；黄色描边：检索与生成部件；红色：交付的回答与核验。
\begin{tikzpicture}[
  x=1mm,y=1mm,
  figure picture,
  inp/.style={figure node,draw=FigureInput,fill=FigureInputFill,
    minimum height=11mm,minimum width=20mm,inner xsep=2pt,inner ysep=2pt,
    font=\figurefont\scriptsize},
  mod/.style={figure node,draw=FigureModel,fill=FigureModelFill,
    minimum height=11mm,minimum width=21mm,inner xsep=2pt,inner ysep=2pt,
    font=\figurefont\scriptsize},
  res/.style={figure node,draw=FigureOutput,fill=FigureOutputFill,
    minimum height=11mm,minimum width=21mm,inner xsep=2pt,inner ysep=2pt,
    font=\figurefont\scriptsize},
  note/.style={font=\figurefont\scriptsize,text=FigureMuted,align=center},
  fl/.style={figure flow},
  back/.style={figure flow,draw=FigureMuted,dash pattern=on 2.5pt off 1.5pt}
]
  \node[inp] (q) at (10,20) {问题$q$};
  \node[mod] (rw) at (36,20) {查询构造\\$q\to\tilde q$};
  \node[mod] (rt) at (62,20) {初检\\BM25或DPR\\前$K$项};
  \node[mod] (rr) at (88,20) {重排\\交叉编码器\\前$k$项};
  \node[mod] (gn) at (114,20) {生成器\\RAG或FiD};
  \node[res] (ct) at (140,20) {带引用回答\\逐陈述核验};

  \draw[fl] (q) -- (rw);
  \draw[fl] (rw) -- (rt);
  \draw[fl] (rt) -- (rr);
  \draw[fl] (rr) -- (gn);
  \draw[fl] (gn) -- (ct);

  % 再检索回路
  \draw[back] (gn.north) -- ++(0,7) -| node[pos=.25,above,note] {Self-RAG：反思标记决定是否检索，并评判片段与生成} (rt.north);
  \draw[back] ([xshift=6mm]gn.north) -- ++(0,15) -| node[pos=.25,above,note] {IRCoT：上一句推理作为新查询，片段并入上下文} ([xshift=-6mm]rt.north);

  % 方法标注
  \node[note] at (36,7) {查询改写\\HyDE假设文档};
  \node[note] at (62,7) {检索器训练\\见第11章};
  \node[note] at (88,7) {以计算换\\排序精度};
  \node[note] at (114,7) {片段边缘化\\或解码器融合};
  \node[note] at (140,7) {ALCE引用\\召回与精确率};
\end{tikzpicture}
